\documentclass[12pt]{article}

\usepackage{amsmath}
\usepackage{amsfonts}
\usepackage{amssymb}
\usepackage{mathrsfs}
\usepackage{physics}
\author{Yuchen Qian}
\date{\today}
\title{Symmetry in Exact diagonalization method}

\begin{document}
\maketitle
The central problem of many-body physics is to diagonalize the Hamiltonian $\mathcal{H}$ in the Hilbert space $\mathscr{H}$. To be explicit, in computational basis we have\begin{align}
    \mathscr{H} = \mathrm{span}\{\ket{\sigma_1\dots\sigma_L}\mid \sigma_i=1,\dots,d\}.
\end{align}The problem with a naive approach is that $\mathscr{H}$ is exponentially large. Let $L$ be the lattice size and $d$ be the local Hilbert space dimension, the dimension of $\mathscr{H}$ is $d^L$. If the system has $\mathsf{U}(1)$ symmetry, then $\mathscr{H}$ is further decomposed into subspaces $\mathscr{H}_N$ with fixed particle number $N$, defined as\begin{align}
    \mathscr{H}_N = \mathrm{span}\{\ket{\sigma_1\dots\sigma_L} \mid \sum_{i=1}^L \sigma_i = N\},
\end{align} and the dimension of $\mathscr{H}_N$ is $\binom{L}{N}$, which is still exponentially large. 

To address such problem, we will use the representation theory\cite{fulton1991representation} to decompose the Hilbert space into irreducible representations of the symmetry group $G$ of the Hamiltonian $\mathcal{H}$. 

We will assume $G$ is a finite group, and it acts on the Hilbert space $\mathscr{H}$ via a unitary representation $U: G \to U(\mathscr{H})$. What we mean by $\mathcal{H}$ has symmetry $G$ is that $\mathcal{H}$ commutes with the action of $G$, i.e. $[U(g),\mathcal{H}] = 0$ for all $g\in G$. By Schur's lemma, we know that $\mathcal{H}$ is block diagonal in the basis of irreducible representations of $G$. Hence we can diagonalize $\mathcal{H}$ in each irreducible representation separately, which reduces the computational cost significantly.
\section{Building the symmetry group}
In our case case, we consider two kinds of symmetry: the lattice symmetry that permutes the sites and the internal symmetry that performs a spin flip:\begin{align}
    g\ket{\sigma_1\dots\sigma_L} &= \ket{\sigma_{g^{-1}(1)}\dots\sigma_{g^{-1}(L)}},\quad g\in G_{\mathrm{latt}}\\
    \mathcal{F}\ket{\sigma_1\dots\sigma_L} &= \ket{(d-\sigma_1),\dots ,( d-\sigma_L)}
\end{align}
and the full symmetry group is $G = G_{\mathrm{latt}}\times \mathbb{Z}_2$. The lattice symmetry group $G_{\mathrm{latt}}$ is further decomposed into a semi-direct product of the translation  group $G_{\mathrm{tr}}$ and the point group $G_{\mathrm{pt}}$, i.e. $G_{\mathrm{latt}} = G_{\mathrm{tr}}\rtimes G_{\mathrm{pt}}$, where $G_{\mathrm{tr}} = \mathbb{Z}_{N_1}\times\cdots\mathbb{Z}_{N_d}$ is a finite Abelian group and $G_{\mathrm{pt}}$ is a finite group. 

Suppose we already know all the information of $G_{\mathrm{pt}}$, we need to \begin{enumerate}
    \item Find the conjugacy classes of $G_{\mathrm{latt}}$ and $G$.
    \item Calculate the character table of $G_{\mathrm{latt}}$ and $G$.
\end{enumerate}
\subsection{The lattice Group}
First we describe the structure of the lattice group $G_{\mathrm{latt}}=G_{\mathrm{tr}}\rtimes G_{\mathrm{pt}}$. Let $\mathbf{t}\in G_{\mathrm{tr}} $ and $r\in G_{\mathrm{pt}}$. The element in $G_{\mathrm{latt}}$ can be written as $(\mathbf{t},r)$. 
It acts on the real space as\begin{align}
    (\mathbf{t},r)\mathbf{x} = r\mathbf{x} + \mathbf{t}.
\end{align}The group multiplication defined as\begin{align}
    (\mathbf{t}_1,r_1)(\mathbf{t}_2,r_2) = (\mathbf{t}_1+r_1\mathbf{t}_2,r_1r_2)
\end{align}where $r(\mathbf{t})$ is the action of $r$ on $\mathbf{t}$. The inverse of $(\mathbf{t},r)$ is given by\begin{align}
    (\mathbf{t},r)^{-1} = (-r^{-1}\mathbf{t},r^{-1}).
\end{align}
\paragraph{Conjugacy Class} Since $G_{\mathrm{tr}}$ is a normal subgroup of $G_{\mathrm{latt}}$, the conjugacy class of $G_{\mathrm{latt}}$ is divided into two cases:\begin{enumerate}
    \item 
    Since\begin{align}
        (a,r)(\mathbf{t},1)(a,r)^{-1}&= (a+r\mathbf{t}, r)(-r^{-1}a,r^{-1}) \\
        &= (a+r\mathbf{t}-r^{-1}a,1)
    \end{align}
    for $(\mathbf{t},1)$ and $(\mathbf{t}',1)$ to be in the same conjugacy class, we need $\mathbf{t}' = a+r\mathbf{t}-r^{-1}a$ for some $a\in G_{\mathrm{tr}}$ and $r\in G_{\mathrm{pt}}$.
    \item Translations and point group elements. We have\begin{align}
        (\textbf{a},r)(\mathbf{t},{s})(\textbf{a},r)^{-1} &= (\textbf{a}+r\mathbf{t}, r{s})(-r^{-1}\textbf{a},r^{-1}) \\
        &= (\textbf{a}+r\mathbf{t}-r^{-1}\textbf{a},r{s}r^{-1})
    \end{align}
    Hence for $(\mathbf{t},{s})$ and $(\mathbf{t}',{s}')$ to be in the same conjugacy class, we need $\mathbf{s}$ and $\mathbf{s}'$ to be in the same conjugacy class of $G_{\mathrm{pt}}$, and $\mathbf{t}' = \textbf{a}+r\mathbf{t}-r^{-1}\textbf{a}$ for some $\textbf{a}\in G_{\mathrm{tr}}$ and $r\in G_{\mathrm{pt}}$.
\end{enumerate}
\paragraph{Character} 
The irreducible representations of $G_{\mathrm{tr}}$ are one-dimensional, and are labeled by the momentum \begin{align}
    \textbf{k}\in \mathrm{BZ} = \left\{(k_1,\dots,k_d)\mid k_i = 0,\frac{2\pi}{N_i},\dots,\frac{2\pi(N_i-1)}{N_i}\right\},
\end{align}with the character defined as\begin{align}
    \chi_{\textbf{k}}^{G_{\mathrm{tr}}}(\mathbf{t}) = \textrm{e}^{\textrm{i} \textbf{k}\cdot \textbf{t}}.
\end{align}
The irreducible representations of $G_{\mathrm{latt}}$ are labeled by the momentum $\textbf{k}$ and the irreducible representations of the little group \begin{align}
    G_{\textbf{k}} = \{r\in G_{\mathrm{pt}}\mid r\textbf{k} = \textbf{k}\}\subset G_{\mathrm{pt}}.
\end{align}
Suppose we know all the irreducible representations of the little group $G_{\textbf{k}}$, with its character denoted as $\chi_{\alpha}^{G_{\mathbf{k}}}$, then the irreducible representations of $G_{\mathrm{latt}}$ are given by\begin{align}
    \chi_{\textbf{k},\alpha}^{G_{\mathrm{latt}}}(t,r) = \sum_{g\in G_{\mathrm{pt}}/G_{\textbf{k}}} \chi_{\textbf{k}}^{G_{\mathrm{tr}}}(gt)\chi_{\alpha}^{G_{\mathbf{k}}}(grg^{-1})
\end{align}
where $g$ runs over the left cosets of $G_{\textbf{k}}$ in $G_{\mathrm{pt}}$. 

\bibliography{bib.bib}
\bibliographystyle{plain}

\end{document}